% GPT-6 Astra Pro; /root/fixed_surface; 22 November 2026 UTC.
% Effort not exposed. Separate complete proof for whole independent review.
\section[HD-F1-S38: coupled complexes with adjacent Hom vanishing]{HD-F1-S38: actual flat cohomology extracts the normal obstruction from coupled ambient complexes}
\label{sec:f1-norm-fourfold-adjacent-cohomology-support}

Retain the normal integral norm image $i:Y\hookrightarrow P$,
$E=i_*\mathcal O_Y$, its ideal $I$, and the S29 geometric open.
Write $D=\mathfrak D_P$ and
$W_E=\operatorname{ob}_E(D)$, so $\dim W_E=16-r_C\ge4$.
Take an actual bounded coherent ambient complex $F$ and an index $j$.
Put $C^a=\mathcal H^a(F)$ and require the two ordinary vanishings
\[
 \operatorname{Hom}_P(C^{j-1},C^j)=0,
 \qquad \operatorname{Hom}_P(C^j,C^{j+1})=0.
 \tag{F1S38.1}
\]
The selected actual sheaf $M=C^j$ has set-theoretic support $Y$
and has the S37 generic module
\[
 M_\xi\simeq\bigoplus_a
   (\mathcal O_{P,\xi}/I_\xi^a)^{m_a},\qquad
 \ell_M=\sum_a m_a\binom{a+3}{4}>0.
 \tag{F1S38.2}
\]
All other cohomology, differentials and extension maps of $F$ are
retained. We prove the actual characteristic-image surjection
\[
\begin{split}
 \mathcal T_F:\operatorname{ev}_F
       \bigl(D\oplus H^2(P,\mathcal O_P)\bigr)&\longrightarrow W_E,\\
 \operatorname{ev}_F(\kappa+\eta)&\longmapsto
                                  -\ell_M\operatorname{ob}_E(\kappa).
\end{split}
 \tag{F1S38.3}
\]
For the entire S14 character this gives an invariant full
degree-two trace kernel of dimension at least $16-r_C$ for
every specified coherent finite point lift. No splitting of the
original complex, and no presentation as $Ri_*\mathcal G$, is
required. Paragraphs 4--7 construct an ordinarily indecomposable,
gluable complex with that full character for which no such
presentation exists.

\paragraph{1. The actual Bockstein maps give flat cohomology.}
Suppose $\operatorname{ev}_F(\kappa+\eta)=0$.
The ambient $P$ is smooth, so $F$ is perfect. The complete Toda
construction with the fixed BF/jet sign of S33 gives an actual
perfect scalar-twisted lift $\widetilde F$ on $P_\kappa$.
For the central closed immersion $s:P\hookrightarrow P_\kappa$,
its actual epsilon triangle is
\[
 s_*F\xrightarrow{\,u\,}\widetilde F
       \xrightarrow{\,v\,}s_*F\longrightarrow s_*F[1],
 \qquad vu=0,\qquad uv=\varepsilon\,1_{\widetilde F}.
 \tag{F1S38.4}
\]
On a local finite free representative it is the short exact
sequence of complexes
$0\to\varepsilon L^\bullet\to L^\bullet
\to L^\bullet/\varepsilon L^\bullet\to0$;
multiplication by epsilon identifies the first and last complexes.
Thus (F1S38.4) retains every differential of the actual lift.

The connecting maps on central cohomology are actual sheaf maps
\[
 b_a:C^a\longrightarrow C^{a+1}.
 \tag{F1S38.5}
\]
They are $\mathcal O_P$-linear because both their source and target
are annihilated by epsilon. They glue globally: cohomology of
the actual twisted transition isomorphisms commutes with the
triangle, and their scalar triple product reduces to one on
the central fibre. These are the Bockstein maps of this lift,
not presumed zero maps on a chosen graded replacement.

Both $b_{j-1}$ and $b_j$ vanish by (F1S38.1).
The precise segment of the cohomology sequence therefore gives
\[
 0\longrightarrow s_*M\xrightarrow{u_j}
       \widetilde M:=\mathcal H^j(\widetilde F)
         \xrightarrow{v_j}s_*M\longrightarrow0,
 \qquad u_jv_j=\varepsilon\,1_{\widetilde M}.
 \tag{F1S38.6}
\]
Here $u_j$ is injective and $v_j$ is surjective. It follows that
$\ker\varepsilon=\ker v_j=\operatorname{im}u_j
=\operatorname{im}\varepsilon$, and
$\widetilde M/\varepsilon\widetilde M=M$.
Hence $\widetilde M$ is an actual flat sheaf over the dual
numbers. Its semilinear overlap maps have the same scalar twist
$1-\varepsilon\eta_{ijk}$. No assertion of flatness is made for
the other lifted cohomology sheaves.

\paragraph{2. Applying the actual mixed-power normal map.}
Apply S37 paragraphs 1--4 directly to the flat sheaf
$\widetilde M$ with central fibre (F1S38.2).
Those paragraphs start from such a sheaf and construct the
diagonal ideal-action maps, their length-weighted trace, the
regular conormal values at every defect point, and the actual
glued ideals of a flat deformation of the reduced $Y$.
Thus they give
\[
 \operatorname{ev}_F(\kappa+\eta)=0
       \quad\Longrightarrow\quad
       \eta_Y(\kappa)=0
       \quad\Longrightarrow\quad
       \operatorname{ob}_E(\kappa)=0.
 \tag{F1S38.7}
\]
The first implication uses (F1S38.4)--(F1S38.6) followed by the
actual sheaf trace; the second is the exact S24 map
$\operatorname{ob}_E=d_2^{\mathrm{ideal}}j\eta_Y$.
There is no appeal to an unproved commutation of the Hochschild
action with truncation, or to a converse deformation criterion.

Subtracting two representatives proves that (F1S38.3) is
well defined and linear. It is onto because $\ell_M\ne0$ and
$W_E=\operatorname{ob}_E(D)$. The sign agrees with S31,
$\operatorname{ev}_E(\kappa)=-\operatorname{ob}_E(\kappa)$.
The normalization uses the positive length of the selected
cohomology sheaf, not a presumed splitting or Euler rank of $F$.

\paragraph{3. Whole character and every given point lift.}
Retain the entire equality for the actual complex
\[
 e^B\operatorname{ch}(F)=\lambda t+Q(\beta),\qquad
 B\in H^{1,1}(P,\mathbf Q),\quad
 0\ne\lambda t\in W_K(P),\quad Q\in\mathbf Q[x].
 \tag{F1S38.8}
\]
The same full finite-exponential derivative and BF module map
as in S33--S37 give
\[
 U_F:=\operatorname{ev}_F
   \{\kappa+\kappa\mathbin{\lrcorner}B:\kappa\in D\}
       \subseteq\ker\sigma_F^2,
 \qquad
 -\ell_M^{-1}\mathcal T_F:U_F\twoheadrightarrow W_E.
 \tag{F1S38.9}
\]
Every specified point lift fixes the actual image $U_F$ by
characteristic-action naturality, since translations and
$\operatorname{Pic}^0$ tensoring fix its Hochschild source.
The exact R05 identity-sector and induction/forgetting maps
therefore give the asserted invariant kernel bound. All these
actions are on the original coupled $F$.
The negative-self-Ext requirement is separate for this general
class; it will be verified for the actual example below.

\paragraph{4. Three point ideals and their exact Hom groups.}
Choose distinct closed points $p,q,r$ in $Y_{\mathrm{reg}}$.
For $s$ one of them, put
\[
 J_s=\mathcal I_{s/Y},\qquad E_s=i_*J_s,
 \qquad 0\longrightarrow E_s\xrightarrow{a_s}E
                    \longrightarrow\mathbf C_s\longrightarrow0.
 \tag{F1S38.10}
\]
These are actual sheaves and exact sequences. For any $s,t$,
an ordinary morphism $J_s\to J_t$ is multiplication by a
rational function at the generic point. At every height-one
point both ideals are the structure sheaf, because the chosen
points have codimension four. That rational function is regular
in every height-one local ring. The normal-domain intersection
therefore extends it to a global regular function on $Y$,
which is constant since $Y$ is connected and projective.
For $s\ne t$, a nonzero constant cannot send $J_s$ into $J_t$:
at $t$ the source is $\mathcal O_{Y,t}$ and the target is its
proper maximal ideal. Exact restriction of scalars for the
closed immersion gives
\[
 \operatorname{End}_P(E_s)=\mathbf C,
 \qquad \operatorname{Hom}_P(E_s,E_t)=0\quad(s\ne t).
 \tag{F1S38.11}
\]
In particular, for
\[
             A=E_p\oplus E_r,\qquad T=E_q,
 \tag{F1S38.12}
\]
both cross Hom groups vanish, while
$\operatorname{End}(A)=\mathbf C^2$ and
$\operatorname{End}(T)=\mathbf C$.

\paragraph{5. Actual ambient couplings outside the support Ext image.}
Every $w\in W_E$ is locally zero as an ambient self-Ext class.
Indeed S24 represents it by the actual Cech normal cocycle in
$H^1(Y,\mathscr N)$, maps this by $j$ to
$\operatorname{Ext}^1_P(I,E)$, and then composes with the ideal
boundary. On an affine open the first cohomology class vanishes;
the ideal composition vanishes there as well. This is the
zero local-to-global edge assertion, with its actual maps.

For a smooth point $s$, the object
$R\mathcal Hom_P(E,\mathbf C_s)$ has cohomology supported at
that point; higher cohomology of those sheaves is zero.
Consequently its global Ext groups are its local Ext groups.
The composition
\[
 \operatorname{Ext}^2_P(E,E)\longrightarrow
                        \operatorname{Ext}^2_P(E,\mathbf C_s)
 \tag{F1S38.13}
\]
vanishes on $W_E$, by the preceding local vanishing.
The actual target sequence (F1S38.10) thus lifts every such $w$
to some $\widetilde\delta_s\in\operatorname{Ext}^2_P(E,E_s)$.
No uniqueness of this target lift is asserted.

For $q\ne s$, there is an actual isomorphism
\[
 a_q^*:\operatorname{Ext}^2_P(E,E_s)
                 \xrightarrow{\sim}\operatorname{Ext}^2_P(E_q,E_s).
 \tag{F1S38.14}
\]
Its precise predecessor and successor are
$\operatorname{Ext}^2_P(\mathbf C_q,E_s)$ and
$\operatorname{Ext}^3_P(\mathbf C_q,E_s)$, both zero.
To verify these vanishings, work at $q$, where $E_s=E$ and
$Y$ is smooth. Choose four normal and four support parameters
in the regular local ring of $P$. Apply Hom into $E$ to their
Koszul resolution of the residue field. The normal differentials
become zero; the four support parameters form an $E$-regular
sequence. Their cochain Koszul complex has cohomology only in
degree four. Tensoring the four zero-differential normal factors
therefore has no cohomology below degree four. Faithfully flat
completion may be used for these coordinates, and the finite
Koszul computation descends. Point support again identifies
these local and global Ext groups. This proves (F1S38.14)
through the exact Hom sequence of (F1S38.10).

Choose two nonzero $w_p,w_r\in W_E$; they may be chosen
linearly independent since $\dim W_E\ge4$.
Choose the actual target lifts above and put
\[
 \delta_s=a_q^*(\widetilde\delta_s)
                   \in\operatorname{Ext}^2_P(E_q,E_s),\qquad
 \Psi_{q,s}=(a_s)_*\circ(a_q^*)^{-1},
 \qquad \Psi_{q,s}(\delta_s)=w_s\ne0.
 \tag{F1S38.15}
\]
Both components are nonzero.

Here is the exact comparison with the support category.
On $Y$, the same four support parameters give
$\operatorname{Ext}^2_Y(\mathbf C_q,J_s)=
\operatorname{Ext}^3_Y(\mathbf C_q,J_s)=0$.
The support source sequence and target cohomology sequence
therefore give the actual isomorphisms
\[
\begin{split}
 \operatorname{Ext}^2_Y(J_q,J_s)
    &\xleftarrow[\sim]{a_q^*}
                   \operatorname{Ext}^2_Y(\mathcal O_Y,J_s)
                         =H^2(Y,J_s),\\
 H^2(Y,J_s)&\xrightarrow[\sim]{(a_s)_*}H^2(Y,\mathcal O_Y).
\end{split}
 \tag{F1S38.16}
\]
The last isomorphism uses the vanishing of positive cohomology
of the point sheaf. Exact derived pushforward commutes with
all the displayed arrows. Thus the image of
$\operatorname{Ext}^2_Y(J_q,J_s)$ under ambient pushforward
and then $\Psi_{q,s}$ is precisely the scalar image
\[
 \iota_2:H^2(Y,\mathcal O_Y)\longrightarrow
                            \operatorname{Ext}^2_P(E,E).
 \tag{F1S38.17}
\]
This image has zero intersection with $W_E$.
The actual retraction is precomposition with the quotient unit
$\mathcal O_P\to E$:
\[
 r_2:\operatorname{Ext}^2_P(E,E)\longrightarrow
            \operatorname{Ext}^2_P(\mathcal O_P,E)=H^2(Y,\mathcal O_Y),
 \qquad r_2\iota_2=1,\qquad r_2(W_E)=0.
 \tag{F1S38.18}
\]
The first identity is the scalar unit/multiplication retract of
S18. The second is the S23 ideal-boundary factorization, since
the ideal boundary composed with the unit is zero; equivalently
it is the actual Atiyah naturality at $\mathcal O_P$.
Equations (F1S38.15)--(F1S38.18) show that neither selected
$\delta_s$ lies in the support Ext image. This conclusion uses
the displayed comparison morphism, not a categorical distinction
without a map.

\paragraph{6. A gluable, ordinarily indecomposable ambient complex.}
Form the actual distinguished triangle with
$\delta=(\delta_p,\delta_r):T\to A[2]$:
\[
 A\xrightarrow{i_A}F\xrightarrow{p_T}T[-1]
                       \xrightarrow{\delta[-1]}A[1].
 \tag{F1S38.19}
\]
Its cohomology is $C^0=A$, $C^1=T$ and zero otherwise.
For $j=0$, conditions (F1S38.1) hold by (F1S38.11), and
(F1S38.2) is the actual generic module
$A_\xi=\mathbf C(Y)^2$, of length two.

Negative self-Ext vanishes. For $n\le-2$, the actual
t-structure gives
$\operatorname{Hom}(F,F[n])=0$, because
$F\in D^{\le1}$ and $F[n]\in D^{\ge2}$.
For degree minus one, the two exact Hom sequences of
(F1S38.19), with their negative sheaf-Ext terms zero, give
\[
 \operatorname{Hom}_P(T,A)\xrightarrow{\sim}
              \operatorname{Hom}_P(F,F[-1]),\qquad
 h\longmapsto i_A[-1]h[-1]p_T.
 \tag{F1S38.20}
\]
Explicitly precomposition with $p_T$ identifies the latter
group with $\operatorname{Hom}(T,F)$ after shifting;
postcomposition with $i_A$ identifies that group with
$\operatorname{Hom}(T,A)$. The preceding groups are negative
Ext groups between the sheaves $A,T$ and are zero.
Equation (F1S38.11) proves gluability.

Every ordinary endomorphism $e$ of $F$ acts on its two
cohomology sheaves. Naturality of the truncation triangle gives
\[
 (a_p,a_r,t)\in\mathbf C^2\oplus\mathbf C,
       \qquad a_p\delta_p=t\delta_p,\quad
                       a_r\delta_r=t\delta_r.
 \tag{F1S38.21}
\]
Both components of $\delta$ are nonzero, so $a_p=a_r=t$.
The kernel of this endomorphism map is square-zero.
Indeed, if $e$ has zero cohomology action, then
$ei_A=0$ and $p_Te=0$. The exact Hom sequences give
$e=h p_T$, then $p_T h=0$, and finally
\[
                    e=i_A v p_T,
                 \qquad v\in\operatorname{Ext}^1_P(T,A).
 \tag{F1S38.22}
\]
For the first two zero equalities one uses the actual
isomorphisms $\operatorname{Hom}(A,F)=\operatorname{End}(A)$
and $\operatorname{Hom}(F,T[-1])=\operatorname{End}(T)$.
For $p_T h=0$, precomposition with $p_T$ on
$\operatorname{End}(T[-1])$ is injective, its preceding
group being $\operatorname{Hom}(A[1],T[-1])=0$.
The final factorization is the target Hom sequence of the
same triangle. Conversely (F1S38.22) kills cohomology.
Products of two such endomorphisms vanish because
$p_Ti_A=0$.

The image of $\operatorname{End}(F)$ is exactly the scalar
diagonal, since scalar identities occur. An idempotent has
scalar image zero or one. In the first case it lies in the
square-zero kernel and is zero; in the second case its
complement is zero. Thus $F$ is ordinarily indecomposable.

It is not isomorphic to $Ri_*\mathcal G$ for any actual
$\mathcal G\in D^b_{\mathrm{coh}}(Y)$.
Such a presentation would identify its two cohomology sheaves
with $J_p\oplus J_r$ and $J_q$, because exact closed pushforward
is faithful on sheaves and commutes with truncations.
Its actual truncation coupling would then be the pushforward
of an element of
$\operatorname{Ext}^2_Y(J_q,J_p\oplus J_r)$.
Changing those cohomology identifications only rescales the
two target summands and the source by (F1S38.11), so it preserves
that image. This contradicts the explicitly computed
non-membership of both components in (F1S38.15)--(F1S38.18).

\paragraph{7. Its entire S14 character and the remaining scope.}
The actual triangle and the three point sequences give, after
the rational cycle-class map,
\[
\begin{split}
 \operatorname{ch}(F)
   &=\operatorname{ch}(E)-[p]-[r]+[q]\\
   &=\operatorname{ch}(E)-\frac{\beta^8}{\int_P\beta^8}.
\end{split}
 \tag{F1S38.23}
\]
Equality of the three point classes is used only in top
rational cohomology; their actual Chow classes are retained
before applying that map. S14 gives
$\operatorname{ch}(E)=w+Q_E(\beta)$, with $w\ne0$.
Thus (F1S38.23) is the complete required character with $B=0$,
the same determinant component, and the displayed change of
the top polynomial coefficient. The complex exists on the
actual norm image, is gluable and ordinarily indecomposable,
and has the invariant trace kernel (F1S38.9). The trivial
point group supplies one actual coherent lift; the theorem
applies to every other specified lift that exists, without
postulating such a lift.

The proof does not claim that arbitrary lifted cohomology is
flat. If an adjacent Bockstein is nonzero, the exact sequence
(F1S38.6) is no longer available. The separate fixed-coefficient
moment calculation for complexes with generic reduced
cohomology does not by itself prove coordinate independence
or the global normal-map descent. Those unrestricted coupled
complexes remain outside this result.

\paragraph{Primary-source scope.}
The actual scalar-twisted perfect lift, local finite free
representatives and signs retain the fully read Toda
Proposition 6.1 and Lemmas 6.2--6.4 scopes of S33--S34.
The epsilon sequence and every new Bockstein/flatness map
are exhibited in (F1S38.4)--(F1S38.6).
The mixed-power ideal map is the complete reviewed S37
construction. Exact Hom sequences and triangle factorizations
use the complete Stacks 0149 proof and 05QK definitions already
read for R33--R34. The normal height-one intersection is the
fully read Stacks 0AVB scope of S33. The required point Ext
vanishings are computed directly by the actual finite Koszul
complexes above; no global regularity of singular $Y$ is assumed.
The original local-to-global, scalar retraction and normal
obstruction are the exact S18/S23/S24 maps, reread at their
displayed scopes. No existence statement for a replacement
object or an unspecified derived lift is used.
